% Esquisse d'un programme (cote 119, copy 2): archive cote 119, site page 96 (PDF page 97), typescript page 18.
% typescript. Draft compared: esquisse/pages/esq_119_020.tex (copy 1, site 20).
% AI-transcribed by Claude Opus 5.5 (claude-opus-5-5, adaptive thinking, effort high), Message Batches API, batch esq119_full (msgbatch_015xqDHkTumDysRF1qngZhRJ), 2026-11-13; not yet checked by eye.
\scanpage{119/96}
\mspage{18}

Il en est qui, face à cela, se contentent de hausser les
épaules d'un air  désabusé et de parier qu'il n'y a rien à tirer
de tout cela, sauf des rêves. Ils oublient, ou ignorent, que notre
science, et toute science, serait bien peu  de choses, si depuis
ses origines elle n'avait été nourrie des rêves et des visions
de ceux qui s'y adonnent avec passion.

4.        Dès le début de ma réflexion sur les cartes bidimension-%
nelles, je me suis intéressé plus particulièrement aux cartes dites
``\uline{régulières}'', c'est-à-dire celles dont le groupe des automorphismes
opère transitivement (et de ce fait, de façon simplement transitive)
sur l'ensemble des repères. Dans le cas orienté et en termes de
l'interprétation algébrico-géométrique du paragraphe précédent,
ce sont les cartes qui correspondent à un revêtement \uline{galoisien} de
la droite projective. Très vite aussi, et dès avant même qu'appa-%
raisse le lien avec la géométrie algébrique, il apparaît néces-%
saire aussi de ne pas exclure les cartes infinies, qui intervien-%
nent notamment de façon naturelle comme revêtements universels des
cartes finies. Il apparaît (comme conséquence immédiate du ``diction-%
naire'' des cartes, étendu au cas des cartes pas nécessairement finies)
que pour tout couple d'entiers naturels  $p, q$ \hand{$\geq$} 1, il existe à iso-%
morphisme (non unique) près une carte 1-\uline{connexe} et une seule qui
soit de type  $(p,q)$  i.e. dont tous les sommets soient d'ordre  $p$
et toutes les faces d'ordre  $q$, et cette carte est une carte régu-%
lière. Elle se trouve épinglée par le choix d'un repère, et son
groupe des automorphismes est alors canoniquement isomorphe au quo-%
tient du groupe cartographique (resp. du groupe cartographique
orienté, dans le cas orienté) par les relations supplémentaires
\[ \rho_s^{\,p} = \rho_f^{\,q} = 1. \]
Le cas où ce groupe est fini est le cas ``pythagoricien'' des cartes
régulières sphériques, le cas où il est infini donne les pavages

\margin{96}
